% =====================================================================
% THÔNG TIN BÁO CÁO — chỉ cần sửa file này cho mỗi báo cáo mới.
% Các giá trị được dùng ở trang bìa, header và footer.
% =====================================================================

% --- Đơn vị ----------------------------------------------------------
\newcommand{\TenDaiHoc}{ĐẠI HỌC QUỐC GIA THÀNH PHỐ HỒ CHÍ MINH}
\newcommand{\TenTruong}{TRƯỜNG ĐẠI HỌC BÁCH KHOA}
\newcommand{\TenTruongHeader}{TRƯỜNG ĐẠI HỌC BÁCH KHOA - ĐHQG-HCM}
\newcommand{\TenKhoa}{KHOA KHOA HỌC VÀ KỸ THUẬT MÁY TÍNH}
% Dạng viết thường, dùng giữa câu văn (lời cảm ơn, cam đoan)
\newcommand{\TenTruongThuong}{Trường Đại học Bách khoa -- ĐHQG-HCM}
\newcommand{\TenKhoaThuong}{Khoa Khoa học và Kỹ thuật Máy tính}

% --- Báo cáo ---------------------------------------------------------
\newcommand{\LoaiBaoCao}{ĐỒ ÁN CHUYÊN NGÀNH}   % ví dụ: BÁO CÁO BÀI TẬP LỚN
\newcommand{\TenDeTai}{TÊN ĐỀ TÀI CỦA BẠN}
% \newcommand{\TenDeTaiEn}{ENGLISH TITLE}      % bỏ comment và bật trong partials/title.tex nếu cần
\newcommand{\NienKhoa}{Niên khóa 2026-2027}

% --- Người thực hiện -------------------------------------------------
\newcommand{\GiangVien}{TS. Họ Tên Giảng Viên}
\newcommand{\NhomSV}{% mỗi sinh viên một dòng, kết thúc bằng \\
  Nguyễn Văn A -- 2300001 \\
  Trần Thị B -- 2300002 \\
  Lê Văn C -- 2300003 \\
}

% Tên gọi dùng trong lời cảm ơn / cam đoan ("Nhóm chúng em", "Em", ...)
\newcommand{\XungHo}{Nhóm chúng em}
